\documentclass[11pt]{article}
\usepackage[margin=1in]{geometry}
\usepackage[T1]{fontenc}
\usepackage{booktabs,longtable,array,enumitem}
\usepackage[hidelinks]{hyperref}
\usepackage{xurl}
\setlist{itemsep=2pt,topsep=3pt}
\title{Data Request to NJ TRANSIT\\\large River LINE Schedule Optimization, ISE Senior Design Team 22}
\author{Virginia Tech, Grado Department of Industrial and Systems Engineering}
\date{Prepared September 28, 2026}
\begin{document}
\maketitle

\noindent The team has already collected what is publicly available: the September 2026 GTFS schedule, NOAA weather, National Transit Database aggregates, OpenStreetMap track data, and the 2010 DVRPC onboard survey (see the companion document \emph{Data Sources and Provenance}). The items below cannot be obtained publicly.

They are grouped by priority. For each item, the list states what the team would use it for and the smallest useful form. Any format is fine: CSV, Excel, scanned paper, or read-only access to a report.

\section*{Priority 1 --- needed to calibrate the model}
\begin{longtable}{@{}p{0.5cm}p{5.3cm}p{5.2cm}p{3.6cm}@{}}
\toprule
\# & Data & Used for & Minimum useful form \\
\midrule\endhead
1 & \textbf{Hold sheets} since Feb.~14, 2026 (6\,am and 6\,pm): cars available, cars down, and the reason, by car number & The distribution of cars available by day and season, which is the main random input to every policy & Scans or photos of the sheets; the team will transcribe them \\
2 & \textbf{Failure and work-order log} by car: date and time, location, failure type, mileage, date returned to service & Failure rates per car-mile and by season and temperature; repair-time distributions; which cars fail most & Export from the maintenance system (Kinkisharyo), or a spreadsheet \\
3 & \textbf{Train-level on-time and cancellation records} with causes, 2025--2026 & Validating the simulation's on-time and cancellation predictions; separating equipment causes from other causes & Daily CSV by train \\
4 & \textbf{Rescue and recovery practice}: who performs rescues, from where, typical time to clear a disabled car, and whether a 2-car consist can move its partner & The rescue model, where each failure also takes a second car out of service & A short interview with operations is enough \\
\bottomrule
\end{longtable}

\section*{Priority 2 --- needed to make the timetables realistic}
\begin{longtable}{@{}p{0.5cm}p{5.3cm}p{5.2cm}p{3.6cm}@{}}
\toprule
\# & Data & Used for & Minimum useful form \\
\midrule\endhead
5 & \textbf{Track chart}: siding limits, switch locations, signal blocks, and the current speed restrictions and slow orders & The single-track timetable model. Public sources disagree about how much of the line is single track: 12.7 miles per NJDOT, 15.2 miles per OpenStreetMap & PDF track chart and the slow-order bulletin \\
6 & \textbf{Minimum turnback and layover times} at Entertainment Center and Trenton; meet and dwell rules at sidings & Model constraints; the team currently assumes 5 minutes & Operating rule or bulletin \\
7 & \textbf{Crew rules and availability}: operator shifts, extra-board size, relief points & Checking whether operators, rather than cars, limit service on some days & Summary numbers \\
8 & \textbf{Consist rules}: when 2-car trains run and why & Capacity on event days; rescue options & Operating rule \\
9 & \textbf{Name of the scheduling system} (for example HASTUS) and its import format & Delivering tiered schedules in a form NJT can load & System name and a sample export \\
\bottomrule
\end{longtable}

\section*{Priority 3 --- needed for rider impact}
\begin{longtable}{@{}p{0.5cm}p{5.3cm}p{5.2cm}p{3.6cm}@{}}
\toprule
\# & Data & Used for & Minimum useful form \\
\midrule\endhead
10 & \textbf{Ridership by trip or by hour}: automatic passenger counts if the cars have them; otherwise ticket vending machine and mobile-app sales and activations by station and hour & Time-of-day demand for the waiting-time and crowding estimates; the model now assumes an hourly profile & Hourly counts by station for a few typical weeks \\
11 & \textbf{Any recent onboard or customer survey} of River LINE riders (or the raw 2010 DVRPC file) & Origin--destination patterns and transfers & Survey report or file \\
12 & \textbf{Special-event service history}: extra trips for waterfront concerts or ballgames, with ridership & Event planning within the 10\,pm freight window & List of events and the service run \\
13 & \textbf{Rider complaints} tagged River LINE, by month & Relating complaints to cancellations & Monthly counts \\
\bottomrule
\end{longtable}

\section*{Priority 4 --- costs}
\begin{longtable}{@{}p{0.5cm}p{5.3cm}p{5.2cm}p{3.6cm}@{}}
\toprule
\# & Data & Used for & Minimum useful form \\
\midrule\endhead
14 & \textbf{Marginal cost per train-hour} (operator wages, fuel, maintenance per car-mile) & The operating-cost side of the trade-off; the model now brackets it with National Transit Database averages & One or two numbers \\
15 & \textbf{Cost and lead time of supplemental buses} & Bus-bridge alternatives on low-car days & Rate per bus-hour \\
16 & \textbf{Rehabilitation and overhaul schedule} (14-car program) and the fleet replacement timeline & The expected number of cars available over the next two years & Planned dates \\
\bottomrule
\end{longtable}

\section*{Questions for the next client meeting}
\begin{enumerate}
\item Does ``8-car schedule'' mean 8 cars (single-car trains) in service? The September 2026 GTFS shows 8 vehicle blocks active in the peaks and 6 midday.
\item What does ``10 cars a day is the goal'' mean: 10 available, or 8 in service plus 2 protect cars?
\item Where are protect cars kept today, and could one be staged at Burlington or Trenton?
\item How is on-time performance counted: are cancelled trips included?
\item Could a reduced schedule be published by 6\,am (app, GTFS-realtime, station signs) when the hold sheet is short?
\end{enumerate}

\section*{Data handling}
The team will use the data only for this course project, keep it in a private repository, and share results with NJ TRANSIT before any presentation. Contact: Prof.~Robert Hildebrand (faculty advisor).
\end{document}
